% ==============================================================================
% SECTION 5: SECURITY HARDENING, QOS & AUTOMATION STRATEGY (STEPS 7 - 9)
% ==============================================================================
\section{Security Hardening, QoS \& Programmable Automation (Steps 7--9)}

\subsection{Edge NAT/PAT, ACL Micro-Segmentation \& Layer 2 Defense (Step 7)}
\paragraph{Zero Trust Access Control \& NAT Exemption.}
Every campus edge router performs \textbf{PAT/NAT Overload} on \texttt{Gi0/0} (\texttt{ip nat outside}) so internal users access Internet services directly. Crucially, the NAT access list explicitly denies inter-campus traffic (\texttt{10.10.0.0/16}, \texttt{10.20.0.0/16}, \texttt{10.30.0.0/16}, \texttt{192.168.0.0/16}) before permitting Internet translation, preventing NAT from breaking \texttt{Tunnel12}/\texttt{Tunnel13} private routing. At Hanoi, an inbound Extended ACL on \texttt{Gi0/1.50} enforces Zero Trust isolation for \texttt{GUEST\_WIFI} (\texttt{10.20.50.0/25}):

\begin{lstlisting}[style=CiscoCLI, caption={Extended ACL Isolating Hanoi Guest Wi-Fi (VLAN 50) from Corporate Networks.}]
ip access-list extended ACL_GUEST_ISOLATION
 remark Allow DHCP/DNS relay reachability if needed, block all RFC1918 Corporate LANs
 deny   ip 10.20.50.0 0.0.0.127 10.0.0.0 0.255.255.255
 deny   ip 10.20.50.0 0.0.0.127 172.16.0.0 0.15.255.255
 deny   ip 10.20.50.0 0.0.0.127 192.168.0.0 0.0.255.255
 permit ip 10.20.50.0 0.0.0.127 any
!
interface GigabitEthernet0/1.50
 ip access-group ACL_GUEST_ISOLATION in
\end{lstlisting}

\paragraph{Layer 2 Defense-in-Depth \& Device Hardening.}
Access switches enforce \textbf{Port Security} (maximum~2 sticky MAC addresses, \texttt{violation restrict}) to block unauthorized switches and MAC flooding, \textbf{DHCP Snooping} (trusting only trunk uplinks \texttt{Gi0/1--2}) to prevent rogue DHCP servers, and \textbf{Dynamic ARP Inspection (DAI)} to defeat ARP spoofing. Telnet and HTTP are disabled across all 10 devices in favor of \textbf{SSHv2} (\texttt{ip domain name tvi.lab}, 2048-bit RSA keys, \texttt{transport input ssh}).

\subsection{Quality of Service (QoS) for VoIP \& Video Convergence (Step 8)}
Under congestion on WAN uplinks caused by AI dataset synchronization in Hanoi or warehouse telemetry in Da Nang, real-time VoIP calls suffer jitter and packet loss. To guarantee voice clarity, \textbf{Low Latency Queuing (LLQ)} is attached outbound on WAN edge interfaces (\texttt{Gi0/0} and Tunnels), strictly prioritizing VoIP packets marked \texttt{DSCP EF} (Expedited Forwarding, 30\% priority queue) and allocating guaranteed bandwidth to Video Conferencing marked \texttt{DSCP AF31} (20\% CBWFQ), while non-critical data shares the remaining bandwidth via fair queuing.

\begin{lstlisting}[style=CiscoCLI, caption={Standardized LLQ QoS Policy-Map for WAN Interfaces.}]
class-map match-any CM-VOIP
 match dscp ef
class-map match-any CM-VIDEO
 match dscp af31
!
policy-map PM-WAN-QOS
 class CM-VOIP
  priority percent 30
 class CM-VIDEO
  bandwidth percent 20
 class class-default
  fair-queue
\end{lstlisting}

\subsection{Programmable Network Automation Workflow (Step 9)}
To eliminate configuration drift across regional switches, Step~9 implements a modular Python Infrastructure-as-Code workflow located in \texttt{Automation/} using \textbf{Netmiko}, \textbf{PyYAML}, and Python's \texttt{getpass} library:
\begin{itemize}[leftmargin=1.5em, itemsep=2pt]
    \item \textbf{Externalized Inventory (\texttt{Automation/inventory.yaml}):} Stores target switch hostnames, \texttt{VLAN 99} management IPs (\texttt{192.168.1.10}, \texttt{192.168.2.10}, \texttt{192.168.3.10}), and \texttt{cisco\_ios} device profiles outside script logic.
    \item \textbf{Zero Hardcoded Secrets:} SSH credentials are prompted interactively at runtime via \texttt{getpass()}, ensuring zero plaintext secrets in source control.
    \item \textbf{Idempotent Execution, Backup \& Exception Handling:} Before and after applying baseline commands (\texttt{ntp server 10.10.10.50}, MOTD banner, and corporate VLANs \texttt{10, 20, 30, 99}), the script captures \texttt{show running-config} and \texttt{show ip interface brief} into timestamped files inside \texttt{Automation/backups/}. Explicit \texttt{try/except} blocks handle SSH timeout and authentication exceptions so unreachable nodes are recorded in \texttt{automation.log} (inside \texttt{Automation/logs/}) without crashing the batch run.
\end{itemize}

\begin{figure}[H]
\centering
\begin{tikzpicture}[
    node distance=0.85cm and 0.9cm,
    box/.style={draw=PrimaryBlue, fill=SoftBlue, rounded corners=4pt, text width=3.1cm, minimum height=1.15cm, align=center, font=\scriptsize},
    err/.style={draw=WarnAmber, fill=SoftAmber, rounded corners=4pt, text width=3.1cm, minimum height=1.15cm, align=center, font=\scriptsize},
    ok/.style={draw=SuccessGreen, fill=SoftGreen, rounded corners=4pt, text width=3.1cm, minimum height=1.15cm, align=center, font=\scriptsize},
    arr/.style={-{Stealth[length=2.5mm]}, thick, color=SlateDark}
]
\node[box] (inv) {\textbf{1. Load Inventory}\\Read \texttt{inventory.yaml} \& prompt \texttt{getpass()}};
\node[box, right=of inv] (ssh) {\textbf{2. SSHv2 Connect}\\Netmiko \texttt{ConnectHandler} to \texttt{VLAN 99} IP};
\node[ok, right=of ssh] (bak) {\textbf{3. Pre/Post State}\\Save running-config \& interfaces to \texttt{backups/}};
\node[ok, below=of bak] (push) {\textbf{4. Push Baseline}\\Deploy NTP \texttt{.50}, MOTD Banner \& VLANs \texttt{10,20,30,99}};
\node[err, below=of ssh] (exc) {\textbf{Exception Handler}\\Catch Timeout / Auth Fail $\rightarrow$ log to \texttt{automation.log}};

\draw[arr] (inv) -- (ssh);
\draw[arr] (ssh) -- node[above, font=\tiny]{Connected} (bak);
\draw[arr] (bak) -- (push);
\draw[arr] (ssh) -- node[right, font=\tiny]{Error} (exc);
\end{tikzpicture}
\caption{Execution pipeline and fault-isolation logic of \texttt{Automation/auto\_provision.py}.}
\label{fig:automation_pipeline}
\end{figure}

\begin{lstlisting}[style=YamlStyle, caption={Externalized Switch Inventory (\texttt{Automation/inventory.yaml}).}]
devices:
  - hostname: HCMC-CORE-SW01
    ip: 192.168.1.10
    device_type: cisco_ios
  - hostname: HANOI-RND-SW01
    ip: 192.168.2.10
    device_type: cisco_ios
  - hostname: DANANG-LOG-SW01
    ip: 192.168.3.10
    device_type: cisco_ios
\end{lstlisting}
